\section{Bilan et perspectives}

\begin{frame}{État des résultats}
\centering\small
\begin{tabular}{lp{8cm}}
    \textcolor{navy}{\textbf{Domaine}} & \textcolor{navy}{\textbf{Statut}} \\
    \hline
    Alignement modèle-cerveau & Robuste (fMRI, ECoG, MEG), mécanisme débattu \\
    Comportement & Partiel : temps de lecture et N400, avec des limites \\
    Mécanismes & Exemples précis, peu généralisés \\
    Supervision cérébrale & Preuves de principe \\
\end{tabular}

\vspace{0.4em}
{\footnotesize Statuts : estimation de l'auteur d'après les exemples présentés.}
\end{frame}

\begin{frame}{Limites}
\begin{itemize}
    \item \textbf{Corrélation n'est pas mécanisme} : un score d'alignement élevé n'établit pas l'identité du calcul
    \item Des scores élevés peuvent s'expliquer par des facteurs simples (longueur, position dans la phrase)
    \item Données majoritairement en anglais
    \item Données intracrâniennes rares, recueillies chez des patients
\end{itemize}
% TODO : ajouter la référence exacte pour le deuxième point
\end{frame}

\begin{frame}{Conclusions}
\begin{enumerate}
    \item Les modèles de langage \textbf{prédisent} l'activité du réseau du langage
    \item Ils \textbf{expliquent} une partie du comportement humain, mais pas tout
    \item Leurs représentations internes peuvent être \textbf{disséquées} avec les mêmes questions que celles posées au cerveau
\end{enumerate}
\end{frame}

\section*{}
\begin{frame}{}
    \centering
    \Huge\bfseries
    \textcolor{navy}{Merci}\\[0.5em]
    \Large Questions ?
\end{frame}
